\begin{abstract}
一个生成的三维物体可以在外观上足够完整，却仍无法回答“这个部件叫什么、位于哪里”。类别无关分割能给出区域，但这些区域未必对应可解释的部件；人工提供空间提示则增加了处理成本。本文研究面向生成式三维资产的视觉语言引导部件分割与语义标注方法。我们利用 Qwen3-VL 将可见部件描述转化为名称、边界框及正负点构成的结构化记录，驱动冻结 SAM2 生成候选掩码；随后通过有限的候选级视觉反馈修正提示，并以点引导清理和小框优先融合处理掩码重叠，减少小部件在覆盖仲裁中被吞并的风险。融合后的二维引导送入冻结 SegViGen，英文名称与二维部件 ID 则通过独立记录保存。围绕不同结构类型的生成资产，本文建立二维引导、三维输出投影、多视图人工评价及效率可靠性的分层实验协议。\TBD{补入主实验、消融及效率测量后的核心数值和结论；当前表格没有实测结果。}本文关注语言意图到空间分割条件的可追溯转换，不将名称记录等同于已完成的三维表面语义绑定。
\end{abstract}
\noindent\textbf{关键词：}三维部件分割；视觉语言模型；空间提示；SAM2；生成式三维资产；语义标注
